\begin{Solbook}{1}
 \par \noindent  \begin {enumerate} \item Ta có \begin {align*} (3x-2)^2-(x+4)^2&=0\\ \big [(3x-2)-(x+4)\big ]\big [(3x-2)+(x+4)\big ]&=0\\ (2x-6)(4x+2)&=0. \end {align*} Do đó $x=3$ hoặc $x=-\dfrac {1}{2}$. \par \item Nhân cả hai vế với $12>0$, ta được \begin {align*} 3(3x-2)-4(x+5)&<2(5x+1)\\ 5x-26&<10x+2\\ -28&<5x\\ x&>-\dfrac {28}{5}. \end {align*} Vậy bất phương trình có nghiệm $x>-\dfrac {28}{5}$. \par \item Dùng phương pháp cộng đại số, ta có \begin {align*} &\begin {cases} \dfrac {2x-y}{3}+\dfrac {x+2y}{2}=\dfrac {5}{3},\\ 3(x-1)-2(y+2)=1 \end {cases}\\[0.7em] \Leftrightarrow \quad &\begin {cases} 7x+4y=10,\\ 3x-2y=8 \end {cases}\\[0.7em] \Leftrightarrow \quad &\begin {cases} 7x+4y=10,\\ 6x-4y=16 \end {cases}\\[0.7em] \Leftrightarrow \quad &\begin {cases} 13x=26,\\ 3x-2y=8 \end {cases}\\[0.7em] \Leftrightarrow \quad &\begin {cases} x=2,\\ y=-1. \end {cases} \end {align*} Vậy hệ phương trình có nghiệm $(x;y)=(2;-1)$. \end {enumerate}  \hfill \qedEX 
\end{Solbook}
\begin{Solbook}{2}
 \par \noindent  \begin {enumerate} \setlength {\itemsep }{1em} \par \item Ta có \begin {align*} 23+8\sqrt {7} &=16+2\cdot 4\sqrt {7}+7\\ &=\left (4+\sqrt {7}\right )^2,\\ 23-8\sqrt {7} &=16-2\cdot 4\sqrt {7}+7\\ &=\left (4-\sqrt {7}\right )^2,\\ 14-4\sqrt {10} &=10-2\cdot \sqrt {10}\cdot 2+4\\ &=\left (\sqrt {10}-2\right )^2. \end {align*} Vì $4-\sqrt {7}>0$ và $\sqrt {10}-2>0$ nên \begin {align*} \sqrt {23+8\sqrt {7}}+\sqrt {23-8\sqrt {7}} -\sqrt {14-4\sqrt {10}}\\ &=4+\sqrt {7}+4-\sqrt {7}-\left (\sqrt {10}-2\right )\\ &=10-\sqrt {10}. \end {align*} Vậy biểu thức đã cho rút gọn thành $10-\sqrt {10}$. \par \item Ta có $\sqrt {80}=4\sqrt {5}$, do đó \begin {align*} \dfrac {\sqrt {80}-\sqrt {5}}{\sqrt {5}-\sqrt {3}} +\dfrac {3\sqrt {3}}{\sqrt {5}+\sqrt {3}}\\ &=\dfrac {3\sqrt {5}}{\sqrt {5}-\sqrt {3}} +\dfrac {3\sqrt {3}}{\sqrt {5}+\sqrt {3}}\\ &=\dfrac {3\sqrt {5}\left (\sqrt {5}+\sqrt {3}\right )}{5-3} +\dfrac {3\sqrt {3}\left (\sqrt {5}-\sqrt {3}\right )}{5-3}\\ &=\dfrac {15+3\sqrt {15}+3\sqrt {15}-9}{2}\\ &=3+3\sqrt {15}. \end {align*} Vậy biểu thức đã cho rút gọn thành $3+3\sqrt {15}$. \end {enumerate}  \hfill \qedEX 
\end{Solbook}
\begin{Solbook}{3}
 \par \noindent  Gọi $x$ (km/h) và $y$ (h) lần lượt là vận tốc dự định và thời gian dự định. \par Quãng đường AB là 600 km, xe xuất phát lúc 7 giờ tối hôm trước.  \hfill \qedEX 
\end{Solbook}
\begin{Solbook}{4}
 \par \noindent  \immini {\begin {enumerate} \item $A,M,C,K$ thuộc đường tròn tâm $I$ bán kính $IM$. \item $\triangle MHB=\triangle NHB$, suy ra $\widehat {MBH}=\widehat {NBH}$. \item $\triangle KMH=\triangle KNH$, suy ra $KM=NK$ \par $\widehat {KMO}=180^\circ -\widehat {KMC}-\widehat {OMB}=180-\widehat {KCM}-\widehat {OBM}=\widehat {CKB}=90^\circ $ \end {enumerate}} {\begin {tikzpicture}[line cap=round,line join=round,font=\footnotesize ] \def \r {2} \path (-\r ,0) coordinate (A) (\r ,0) coordinate (B) (0,0) coordinate (O) ($(O)!1!-50:(A)$) coordinate (M) ($(A)!(M)!(B)$) coordinate (H) ($(M)!2!(H)$) coordinate (N) ($(B)!2!(M)$) coordinate (Bt) ($(N)!2.5!(A)$) coordinate (Nt) (intersection of N--Nt and B--Bt) coordinate (C) ($(A)!(C)!(B)$) coordinate (K) ($(C)!2!(K)$) coordinate (Ct) ($(B)!1.5!(N)$) coordinate (Nt) (intersection of C--Ct and B--Nt) coordinate (D) ($(C)!0.5!(A)$) coordinate (I) ; \draw (O) circle (\r ) (M)--(N) (K)--(A)--(B)--(C)--(D)--(B) (C)--(N) (M)--(A)--(B)--cycle (M)--(I)--(K) (M)--(O)--(N)--(K)--(M) ; \foreach \x /\g in{A/-135,B/-45,O/-90,M/90,H/-45,N/-90,C/90,K/180,D/-90,I/70} \draw [fill=white] (\x ) circle (1.5pt) +(\g :.35) node{$\x $}; \end {tikzpicture}}  \hfill \qedEX 
\end{Solbook}
\begin{Solbook}{5}
 \par \noindent  $x$ (triệu đồng) là số tiền ban đầu người đó gửi vào ngân hàng \par $0,5\%\left (1+0,5\%\right )\ge 0,5$ \par Người đó cần gửi số tiền ban đầu ít nhất là 100 triệu đồng.  \hfill \qedEX 
\end{Solbook}
